%----------------------------------------------------------------------------------------
%   CHAPTER 4
%----------------------------------------------------------------------------------------

\chapter{Przeprowadzenie badania i omówienie wyników}

Niniejszy rozdział przedstawia metodologię oraz wyniki badania
przeprowadzonego z udziałem trzech uczestników. Celem eksperymentu
była praktyczna weryfikacja działania systemu w warunkach rzeczywistych,
obejmująca ocenę skuteczności klasyfikacji komend, użyteczności interfejsu
oraz subiektywnych wrażeń użytkowników. Badanie przeprowadzono zgodnie
z protokołem stanowiącym Załącznik~\ref{appendix:protokol}.

% -------------------------------------------------------
\section{Protokół badania}

W badaniu wzięły udział trzy osoby, które nie miały wcześniejszego
doświadczenia z interfejsami mózg--komputer. Sesja składała się z sześciu
etapów opisanych szczegółowo w protokole (Załącznik~\ref{appendix:protokol}).

Na początku uczestnik zapoznawał się z celem eksperymentu oraz zasadą działania interfejsu SSVEP. Następnie przeprowadzano montaż elektrod EEG w punktach O1, O2 oraz Oz, kontrolując poziom impedancji w celu zapewnienia poprawnej jakości sygnału. Uczestnik zajmował miejsce w odległości około 60--70~cm od monitora. Po przygotowaniu rozpoczynano krótką sesję zapoznawczą, w trakcie której badany mógł swobodnie eksplorować działanie systemu. W dalszej części realizowano zadanie polegające na sterowaniu kursorem i klikaniu w cele o malejących rozmiarach. Czasy ukończenia zadania z klikaniem rejestrowano ręcznie przy użyciu stopera. Następnie przechodzono do zadania oceniającego skuteczność klasyfikacji komend. Uczestnik wykonywał polecenia wydawane przez prowadzącego w losowej kolejności, a każda próba była klasyfikowana pod względem poprawności odpowiedzi systemu.

Po zakończeniu części eksperymentalnej uczestnik wypełniał ankietę dotyczącą subiektywnej oceny działania interfejsu. Obejmuje ona zarówno odczucia związane z jego użytecznością, jak i poziom obciążenia poznawczego oraz fizycznego (Załącznik~\ref{appendix:ankieta}).

% -------------------------------------------------------
\section{Wyniki badania}
Pierwsze zadanie polegało na przemieszczeniu kursora myszy i kliknięciu
kolejno w trzy czerwone kwadraty o wymiarach $200\times200$, $100\times100$
oraz $50\times50$ pikseli, wyświetlane na ekranie na raz. 
Zadanie powtarzano trzykrotnie. Za każdym razem obiekty znajdowały się w tym samym miejscu. Mierzono łączny czas ukończenia każdego podejścia, obejmujący zarówno czas ruchu kursora,
jak i czas oczekiwania na reakcję klasyfikatora.

Uczestnik oznaczony jako Osoba~3 nie był w stanie ukończyć zadania
w żadnym z podejść z powodu trudności ze skupieniem wzroku
na polach stymulatora przy jednoczesnej obecności wszystkich pięciu bodźców
na ekranie --- szczegółowe omówienie tego przypadku zawiera
sekcja~\ref{sec:dyskusja}.

Wyniki pozostałych dwóch uczestników zestawiono w tabeli~\ref{tab:task1}.

\begin{table}[!ht]
\centering
\caption{Czasy ukończenia zadania sterowania kursorem [s]}
\label{tab:task1}
\begin{tabular}{lccccc}
    \toprule
    \textbf{Uczestnik} & \textbf{Podejście 1} & \textbf{Podejście 2}
        & \textbf{Podejście 3} & \textbf{Średnia} & \textbf{Min / Maks} \\
    \midrule
    Osoba 1 & 261 & 312 & 270 & 281 & 261 / 312 \\
    Osoba 2 & 207 & 310 & 191 & 236 & 191 / 310 \\
    Osoba 3 & \multicolumn{5}{c}{zadanie nieukończone} \\
    \bottomrule
\end{tabular}
\end{table}

Osoba~1 uzyskała stabilne wyniki w pierwszym i trzecim podejściu
(261~s i 270~s), natomiast drugie podejście było wyraźnie dłuższe
(312~s). Osoba~2 osiągnęła najlepszy wynik w trzecim podejściu (191~s),
co sugeruje efekt uczenia się --- użytkownik stopniowo lepiej opanowywał
technikę skupiania wzroku na bodźcu. Znaczące wydłużenie czasu
w drugim podejściu u Osoby~2 (310~s) może wynikać z chwilowego
zmęczenia lub trudności z trafieniem kursorem w mniejsze kwadraty,
co potwierdzają późniejsze odpowiedzi ankietowe tej osoby.

Warto odnotować, że minimalne opóźnienie systemu wynoszące 4,2~s
(wynikające z długości okna analizy i bufora decyzyjnego, zgodnie
z równaniem~\ref{eq:delay}) jest wliczone w każdy zarejestrowany czas.
Przy konieczności wykonania kilkunastu komend na jedno podejście,
samo opóźnienie systemu stanowi istotną składową łącznego czasu zadania.

% -------------------------------------------------------
Drugie zadanie polegało na wykonaniu przez uczestnika pięciu komend
(prawo, lewo, góra, dół, kliknięcie) w trzech podejściach,
po jednej komendzie na bodziec. Prowadzący wydawał polecenia ustnie
w losowej kolejności. Wynik każdej próby
klasyfikowano jako sukces~(S), niepowodzenie~(N) --- brak reakcji systemu
w czasie poniżej 8~s --- lub błąd~(B) --- wykonanie innej komendy
niż oczekiwana. Odnotowywano również przypadki wielokrotnej
aktywacji komendy (SX2, SX3), gdy bufor decyzyjny wyzwalał
tę samą komendę kilkukrotnie podczas jednego skupienia wzroku.

Zbiorcze wyniki skuteczności przedstawia tabela~\ref{tab:task2_summary},
natomiast rozkład wyników per komenda --- tabela~\ref{tab:task2_per_cmd}.

\begin{table}[!ht]
\centering
\caption{Zbiorcza skuteczność klasyfikacji komend (15 prób na uczestnika)}
\label{tab:task2_summary}
\begin{tabular}{lcccc}
    \toprule
    \textbf{Uczestnik} & \textbf{Sukces} & \textbf{Błąd (B)}
        & \textbf{Niepowodzenie (N)} & \textbf{Wielokrotne} \\
    \midrule
    Osoba 1 & 13/15 (87\%) & 0/15 (0\%)  & 2/15 (13\%) & 2/15 \\
    Osoba 2 & 15/15 (100\%) & 0/15 (0\%) & 0/15 (0\%)  & 13/15 \\
    Osoba 3 & 9/15 (60\%)  & 5/15 (33\%) & 1/15 (7\%)  & 4/15 \\
    \bottomrule
\end{tabular}
\end{table}

\begin{table}[!ht]
\centering
\caption{Skuteczność per komenda --- liczba Sukcesów / Błędów / Niepowodzeń
w 3 podejściach}
\label{tab:task2_per_cmd}
\begin{tabular}{lccccc}
    \toprule
    \textbf{Komenda} & \multicolumn{2}{c}{\textbf{Osoba 1}}
        & \multicolumn{2}{c}{\textbf{Osoba 2}}
        & \textbf{Osoba 3} \\
    \cmidrule(lr){2--3}\cmidrule(lr){4--5}\cmidrule(l){6--6}
    & S & B/N & S & B/N & S / B / N \\
    \midrule
    Prawo  (10,0 Hz) & 3/3 & 0 & 3/3 & 0     & 2/3 / 1B / 0N \\
    Lewo   (8,57 Hz) & 3/3 & 0 & 3/3 & 0     & 1/3 / 2B / 0N \\
    Góra   (6,66 Hz) & 1/3 & 2N & 3/3 & 0    & 1/3 / 1B / 1N \\
    Dół    (7,5 Hz)  & 3/3 & 0 & 3/3 & 0     & 3/3 / 0B / 0N \\
    Klik   (12,0 Hz) & 3/3 & 0 & 3/3 & 0     & 2/3 / 1B / 0N \\
    \bottomrule
\end{tabular}
\end{table}

Osoba~1 osiągnęła skuteczność 87\%, bez ani jednego błędnego
zaklasyfikowania komendy. Jedynym problemem były dwa niepowodzenia
przy komendzie góra (częstotliwość 6,66~Hz) --- klasyfikator
nie przekroczył progu $\theta = 0{,}3$ w dopuszczalnym czasie.
Możliwą przyczyną jest niższa czułość uczestnika na tę konkretną częstotliwość stymulacji.

Osoba~2 uzyskała 100\% skuteczności klasyfikacji, jednak aż
w 13 z 15 prób komenda została aktywowana wielokrotnie
zanim użytkownik oderwał wzrok od bodźca. Świadczy to o bardzo
silnej odpowiedzi SSVEP u tej osoby przy jednoczesnym zbyt krótkim
czasie blokady cooldown (0,5~s) w stosunku do czasu reakcji.
Wartość progu $\theta = 0{,}3$ okazała się zbyt niska
dla tego uczestnika --- przy silnym sygnale system reagował
wielokrotnie w czasie jednego skupienia wzroku.

Osoba~3 osiągnęła skuteczność 60\% przy 33\% błędów,
co jest wynikiem wyraźnie odbiegającym od pozostałych uczestników.
Szczegółowe okoliczności tego przypadku omówiono
w sekcji~\ref{sec:dyskusja}.

% -------------------------------------------------------
\section{Wyniki ankiety}


Po zakończeniu sesji każdy uczestnik wypełnił ankietę oceny systemu
złożoną z trzech części: oceny ogólnej (skala 1--5), oceny obciążenia
podczas zadania (skala 1--5) oraz pytań otwartych. Wzór ankiety zamieszczono w sekcji~\ref{appendix:ankieta}. Wyniki liczbowe zestawiono w tabeli~\ref{tab:ankieta}.

\begin{table}[!ht]
\centering
\caption{Wyniki ankiety --- oceny w skali 1--5}
\label{tab:ankieta}
\begin{tabular}{lp{5.5cm}ccc}
    \toprule
    \textbf{Część} & \textbf{Pytanie / Odczucie}
        & \textbf{O1} & \textbf{O2} & \textbf{O3} \\
    \midrule
    \multirow{6}{*}{Ocena ogólna}
        & System reagował zgodnie z intencją.      & 4 & 3 & 2 \\
        & Czułem/am się w kontroli nad kursorem.   & 4 & 4 & 2 \\
        & Czas oczekiwania był akceptowalny.        & 3 & 4 & 3 \\
        & Bodźce nie powodowały dyskomfortu.        & 5 & 2 & 1 \\
        & Obsługa była intuicyjna.                  & 5 & 3 & 4 \\
        & Skorzystałbym/am ponownie.                & 5 & 2 & 4 \\
    \midrule
    \multirow{4}{*}{Obciążenie}
        & Wymagana koncentracja.                    & 5 & 3 & 5 \\
        & Zmęczenie wzroku.                         & 2 & 5 & 4 \\
        & Frustracja z powodu błędów.               & 3 & 5 & 4 \\
        & Ogólne zmęczenie.                         & 1 & 4 & 3 \\
    \bottomrule
\end{tabular}
\end{table}


Osoba~1 oceniła system pozytywnie --- wysokie noty za intuicyjność
i chęć ponownego użycia (5/5), przy umiarkowanej ocenie
akceptowalności czasu oczekiwania (3/5). Koncentracja wymagana
do obsługi oceniona została bardzo wysoko (5/5), natomiast
ogólne zmęczenie było minimalne (1/5), co sugeruje,
że sesja nie była nadmiernie obciążająca fizycznie.
W odpowiedziach otwartych uczestnik wskazał, że komenda góra
sprawiała największą trudność, natomiast kliknięcie działało
najszybciej i najbardziej przewidywalnie.

Osoba~2 zgłosiła wyraźny dyskomfort wzrokowy (5/5 zmęczenie wzroku)
oraz wysoki poziom frustracji spowodowany wielokrotnymi aktywacjami
komend (5/5). Niska ocena braku dyskomfortu bodźców (2/5)
i chęci ponownego użycia (2/5) koreluje bezpośrednio
z problemem wielokrotnych aktywacji, który był dominującym
doświadczeniem tej osoby. Uczestnik zaproponował ciekawe
techniki skupiania wzroku --- obserwowanie wzorków i cieni
wokół kwadratów zamiast patrzenia w ich centrum.
Sugerował też możliwość poruszania się i mówienia
podczas badania, co jest typową potrzebą użytkowników
niezaznajomionych z wymaganiami rejestracji EEG.

Osoba~3 oceniła system bardzo nisko pod względem kontroli
nad kursorem (2/5) i zgodności z intencją (2/5),
co bezpośrednio wynika z trudności opisanych
w sekcji~\ref{sec:dyskusja}. Mimo to intuicyjność
obsługi i chęć ponownego użycia ocenione zostały
stosunkowo wysoko (4/5), co sugeruje, że uczestnik
rozumiał ideę systemu i dostrzegał jego potencjał
pomimo trudności w sesji. Jako sugestię poprawy
wskazał zmniejszenie jasności i kontrastu monitora,
co mogło łagodzić dyskomfort wynikający z intensywnego
migotania bodźców.

% -------------------------------------------------------
\section{Dyskusja wyników}
\label{sec:dyskusja}

Dane z tabeli~\ref{tab:task2_per_cmd} wskazują na wyraźne
różnice w skuteczności klasyfikacji w zależności od częstotliwości
bodźca. Komenda góra (6,66~Hz) generowała największą liczbę
niepowodzeń u Osoby~1, a u Osoby~3 wykazywała zarówno
błędy, jak i niepowodzenia. Komenda dół (7,5~Hz) była jedyną,
którą Osoba~3 wykonała poprawnie we wszystkich trzech podejściach.
Obserwacja ta jest zgodna z wynikami Osoby~3 przy pracy
z pojedynczym stymulatorem, gdzie najsilniejsze odpowiedzi
rejestrowano właśnie dla 7,5~Hz i 8,57~Hz.

Może to wskazywać na indywidualną wrażliwość kory wzrokowej
na poszczególne częstotliwości stymulacji, niezależną od
ogólnej jakości odpowiedzi SSVEP \cite{guger2012}.
W systemach z większą liczbą użytkowników warta rozważenia
byłaby krótka sesja selekcji częstotliwości przed właściwym
badaniem.


Uzyskane wyniki skuteczności dla Osoby~1 (87\%) i Osoby~2 (100\%)
mieszczą się w zakresie typowym dla systemów SSVEP opartych
na CCA. Systemy opisane przez Gao i współpracowników
\cite{gao2014} osiągały dokładność powyżej 95\% przy oknie
analizy 1--2~s. Zastosowane w niniejszej pracy okno 4~s
zwiększa stabilność klasyfikacji kosztem wydłużonego
czasu reakcji, co bezpośrednio przekłada się
na obserwowane czasy ukończenia zadania pierwszego.

Czasy ukończenia zadania z klikaniem (191--312~s) są wynikiem
relatywnie wysokim w porównaniu z tradycyjną myszą,
jednak zgodnym z oczekiwaniami dla systemu BCI
z minimalnym opóźnieniem 4,2~s i brakiem sprzężenia
zwrotnego wizualnego informującego o stanie klasyfikatora.
Wprowadzenie wizualnego wskaźnika wartości $\rho$
w czasie rzeczywistym mogłoby skrócić czas reakcji
użytkownika i zmniejszyć liczbę niepowodzeń.


Wyniki Osoby~2 ujawniły pewne ograniczenie obecnej implementacji
bufora decyzyjnego. Przy silnej odpowiedzi SSVEP wartość $\rho$
utrzymuje się powyżej progu $\theta$ przez cały czas skupienia
wzroku na bodźcu, co powoduje cykliczne wyzwalanie komend
co 0,5~s (czas blokady cooldown). W 87\% prób Osoby~2
komenda była aktywowana więcej niż jeden raz.

Rozwiązaniem tego problemu może być wydłużenie czasu blokady
lub wprowadzenie mechanizmu detekcji oderwania wzroku
od bodźca, np. na podstawie nagłego spadku wartości $\rho$
poniżej pewnego progu. Alternatywnym podejściem
byłoby automatyczne dostosowanie progu $\theta$ do
charakterystyki sygnału konkretnego użytkownika
w sesji kalibracyjnej.

Osoba~3 stanowi przypadek szczególny, wymagający oddzielnego omówienia.
Uczestnik wykazywał znaczące trudności ze skupieniem wzroku na wybranym
bodźcu w obecności pozostałych czterech stymulatorów na ekranie.
Wykluczona została możliwość nieprawidłowego umiejscowienia elektrod
oraz zbyt wysokiej impedancji kontaktu. Obniżenie progu klasyfikacji
z $\theta = 0{,}3$ do $\theta = 0{,}26$ nie przyniosło
znaczącej poprawy skuteczności.

Dopiero po usunięciu pozostałych stymulatorów z ekranu
i pozostawieniu jednego aktywnego bodźca uczestnik był
w stanie wykonać zadanie skuteczności komend, uzyskując
zadowalające wyniki dla częstotliwości 7,5~Hz i 8,57~Hz,
przy których klasyfikator rejestrował wartości $\rho$
sięgające nawet 0,45. Pozostałe częstotliwości generowały
niestabilne lub słabe odpowiedzi.

Zjawisko to jest opisane w literaturze jako
\textit{BCI inefficiency} --- szacuje się, że około 15--25\%
populacji nie jest w stanie efektywnie korzystać z interfejsów
SSVEP niezależnie od jakości sygnału EEG \cite{vialatte2010}.
Może ono wynikać z indywidualnych różnic neuroanatomicznych
w obszarze kory wzrokowej lub z trudności z kontrolą
uwagi wzrokowej w obecności wielu konkurujących bodźców.
Wynik Osoby~3 wskazuje ponadto na potrzebę adaptacyjnego
doboru częstotliwości stymulacji do indywidualnych
cech użytkownika --- element, który mógłby stanowić
kierunek dalszego rozwoju systemu.




